\section{Experiments and Results}
\label{sec:results}

\emph{Drafting status.} This section will be populated only from the audited
section-level tables. The authoritative patient-grouped GBM summary is still to
be generated before inferential wording is finalised.

At matched section-specific peak budgets, reproduced legacy msiPL obtained mean
mSCF1 values of 0.4024 on CAC and 0.4142 on GBM. The centre-only VAE/GMM/IG
pipeline obtained 0.5285 and 0.5178, respectively, while the fixed-context
variant obtained 0.5595 and 0.5111. These are method-package comparisons:
architecture, training, latent modelling, attribution and ranking all differ
from legacy msiPL.

The spatial comparison is collection-dependent. Fixed context exceeded the
centre-only model on all eight CAC sections in the frozen single-seed
comparison, but did not produce a consistent GBM benefit. Larger-window,
zero-context, shuffled-context and learned-attention controls further prevent a
general claim that correct local topology caused the CAC gain. The final result
presentation will show every paired section and group GBM sections by patient.

Matched-count S3PL obtained higher mean CAC mSCF1 than fixed-context IG under
its native evaluated protocol. On GBM, the executable 10-epoch S3PL baseline
was initialization-sensitive at section level: six of eight sections had a
three-seed range of at least 0.10 mSCF1, although collection means were steadier.
Longer GBM108-positive runs caused initially divergent peak sets and scores to
converge, showing that a favourable short run is not a stable baseline.

Three native 160TopL measurements recorded S3PL workflow totals of 151.2,
173.9 and 173.9 seconds. The VAE workflow totals were 135.3, 121.6 and 121.9
seconds and used slightly less sampled memory. These measurements describe the
specified native procedures; they do not equate ten S3PL epochs with one hundred
VAE epochs or establish universal architectural efficiency.

